\section{Conclusion}
\label{sec:conclusion}

Procurement authorities need triage tools that work with minimal data.
The frequent-loser screen meets this need: it requires only win/loss
records, runs in minutes on any procurement database, and flags
environments with systematically higher conditional prices and strong
proximity to known cartel environments. Applied to 4.5 million
tender-items from Brazilian procurement, it achieves AUC~$= 0.94$
against competition-authority convictions and 3.6--7.7\% higher
conditional prices across four estimation approaches. Five supporting
diagnostics are coherent with coordinated cover bidding, strengthening
the rationale for deployment without claiming to identify the
mechanism.

We propose a three-stage enforcement pathway.
(i)~\emph{Screen}---apply the FL rule to participation records; the
computation requires no bid microdata, no enforcement priors, and no
supervised training.
(ii)~\emph{Triage}---cross-reference FL-flagged markets with network
metrics and oversight-capacity indicators to prioritize the
highest-risk cases; the 12.5$\times$ gradient in the FL--price
association between the smallest and largest purchasing units suggests
that triage resources are best directed toward low-oversight
environments.
(iii)~\emph{Investigate}---deploy bid-level forensic tools
\emph{on the triaged subset}. The screen's near-orthogonality with
bid-level screens (correlation 0.06) means the two stages capture
different information, and combining them reduces both false
positives and missed cases.

Two features make the screen administrable. First, the threshold is
robust: varying the IQR multiplier from 1.0$\times$ to 3.0$\times$
and the win-rate cutoff from 0\% to 5\% produces significant
coefficients across all 36 parameter combinations, so the triage
signal does not hinge on a precise calibration choice. Second, the
data requirements are minimal---participation and outcome records
that every procurement system already collects---making the screen
portable to settings where bid-level forensics are infeasible,
including developing-country procurement systems with limited
investigative capacity.

The screen's main limitation is also its boundary: it flags suspicious
environments, not individual firms, and the conditional price
association does not have a causal interpretation. These are features
of a triage tool, not defects. An enforcement agency deploying the
screen would use it to allocate investigative effort, not to
adjudicate liability.

Replication across procurement systems is the most pressing next step.
The screen's economic rationale---that cover bidders must participate
to perform their role---is institution-general, but its empirical
performance outside BEC remains to be tested. Brazil's recent
procurement reform (Lei 14.133/2021), which restructures bidding
modalities and minimum-bidder rules, provides a natural laboratory
for assessing whether the screen adapts to institutional change.
